% ---- Abstract -------------------------------------------------------
\chapter*{Abstract}
\addcontentsline{toc}{chapter}{Abstract}

Large language models (LLMs) increasingly articulate plausible trading strategies, yet
recent benchmarks show that this fluency does not translate into sound decisions: models
introduce look-ahead leakage into their own backtests, fail to adapt across market
regimes, and execute trades inconsistent with their stated reasoning --- the
\emph{knowledge--action gap}. This thesis asks not whether an LLM agent can trade, but
whether we can build a harness that measures, rigorously and by construction, \emph{how
well it reasons} while doing so.

We design an autonomous trading agent that perceives market data, technical indicators,
news sentiment, portfolio state and order execution \emph{only} through Model Context
Protocol (\mcp) tool servers, operating in a \mbox{perceive\,$\rightarrow$\,reason\,$\rightarrow$\,act\,$\rightarrow$\,observe}
loop. The central methodological contribution is an \emph{\mcp-as-time-machine} design: a
single simulation clock exposes \tnow, and every data tool refuses, at read time, to
return any datum dated strictly after \tnow. The agent calls the identical tools in
backtest and in live operation; only the clock advances. Look-ahead bias is therefore
\emph{structurally impossible} rather than avoided by discipline, and this guarantee is
pinned by tests written before the data layer existed (test-driven development) and
exercised by mutation testing.

On top of this substrate we build a regime-segmented evaluation of reasoning quality with
three complementary metrics: \emph{faithfulness} (does the executed trade match what a
rational, policy-aware judge would do given the same evidence?), measured by a blind
LLM judge and anchored to human annotation; \emph{grounding} (are the numbers cited in the
chain-of-thought real outputs of the tools that were actually called?), computed
deterministically; and a planned KellyBench-style \emph{sophistication} rubric. We report
financial performance against shifted, symmetrically-costed Buy-\&-Hold, momentum and
mean-reversion baselines, with bootstrap confidence intervals.

The primary result is the single-agent system; a Portfolio-Manager multi-agent extension
(analyst reports, a Bull/Bear research debate, a risk debate and a synthesising allocator)
is reported as a bonus that tests whether explicit deliberation closes the
exit-discipline gap observed in the single agent.

\bigskip
\noindent\textit{All quantitative results are reported via centralised macros: the
single-agent figures are from the full two-year canonical run, and the multi-agent figures
from a bounded one-month validation that establishes the pipeline's causal correctness rather
than a full track record; see \cref{ch:results}.}
